\chapter{Representaciones aprendidas, geometría y dinámica del entorno}
\label{ch:latent_env}

Una vez definido el criterio de fidelidad del gemelo, es necesario decidir cómo representar
la escena y el campo de radiofrecuencia. Esta decisión afecta la capacidad de edición del
modelo, el costo computacional, la transferencia entre escenas y la cantidad de mediciones
requerida para la calibración. Por ello, las representaciones neuronales se estudian aquí
como aproximaciones estructuradas del modelo físico y no como sustitutos desligados de él.

\section{Por qué representar y no sólo simular}

El trazado de rayos explícito conserva la interpretabilidad, pero puede ser costoso y depende de una geometría detallada. Los modelos neuronales del campo buscan aproximar
\begin{equation}
(\vect p_T,\vect p_R,G_E)\mapsto H
\end{equation}
a partir de mediciones o representaciones geométricas compactas. La cuestión científica no es únicamente qué modelo reduce el error cuadrático, sino qué representación conserva la estructura necesaria para interpolar, extrapolar, editar escenas, adaptarse con pocas mediciones y soportar inversión.

\section{Regímenes de información}

Para comparar familias heterogéneas se fija primero una consulta $\vect u=(\vect p_T,\vect p_R,f,\vect s_R)$ y un conjunto de pares RF de calibración $\mathcal D_{\rm RF}=\{(\vect u_i,H_i)\}$. Los cuatro regímenes indican información adicional disponible en inferencia:
\begin{align}
\mathcal R_0 &: \{\vect u,\mathcal D_{\rm RF}\},\\
\mathcal R_1 &: \{\vect u,\mathcal D_{\rm RF},G_E^{\rm geom}\},\\
\mathcal R_2 &: \{\vect u,\mathcal D_{\rm RF},G_E^{\rm geom},\{O_j\}\},\\
\mathcal R_3 &: \{\vect u,\mathcal D_{\rm RF},G_E^{\rm geom},\{O_j\},
\vect\theta_{\rm mat},\vect\theta_{\rm ant}\}.
\end{align}
Aquí $G_E^{\rm geom}$ contiene sólo geometría, a diferencia del estado
completo $G_E$ que también incluye materiales. $O_j$ es un objeto segmentado con geometría y clase semántica. El canal reservado para evaluación es siempre la variable que debe predecirse y no una entrada de la consulta. El número de muestras, la región espacial y los metadatos se mantendrán constantes dentro de cada comparación; así se evita presentar como ganancia arquitectónica una ventaja debida a la información disponible.

\section{Campos implícitos}

Un campo neural continuo puede aproximarse como
\begin{equation}
\widehat H=F_\theta(\vect p_T,\vect p_R,f).
\end{equation}
La ventaja es una representación continua y compacta; el riesgo es que la geometría y los materiales queden absorbidos en los pesos, lo que dificulta la edición y la transferencia. NeRF$^2$ representa una línea importante en esta dirección y se adopta como referencia metodológica para campos implícitos \citep{Zhao2023NeRF2}.

\section{Representaciones centradas en objetos e informadas por la geometría}

Una alternativa factoriza el entorno:
\begin{equation}
Z_E=\left\{z_{\rm global},z_{O_1},\ldots,z_{O_M}\right\},
\end{equation}
con
\begin{equation}
z_{O_i}=\left[z_{\rm geom}^{(i)},z_{\rm mat}^{(i)},z_{\rm int}^{(i)}\right].
\end{equation}
Learnable WDT y trabajos condicionados a la geometría como RadTwin constituyen referencias para esta familia \citep{Jiang2025LearnableWDT,Zhang2026RadTwin}. RFScape constituye otro antecedente directo de objetos neuronales acoplados
a trazado editable \citep{Chen2025RFScape}. La hipótesis es que una representación explícitamente estructurada puede necesitar más información geométrica inicial, pero generalizar mejor a cambios de objetos y soportar intervenciones físicas interpretables.

\section{El espacio de trayectorias como representación intermedia}

Se propone representar cada trayectoria mediante un descriptor
\begin{equation}
\vect p_{m,\ell}=
[\Re\alpha_{m,\ell},\Im\alpha_{m,\ell},\tau_{m,\ell},
\theta^{\rm AoA}_{m,\ell},\theta^{\rm AoD}_{m,\ell},
\nu_{m,\ell},\vect c_{m,\ell}],
\end{equation}
donde $\vect c_{m,\ell}$ codifica interacciones; los ángulos se representarán mediante seno/coseno para evitar discontinuidad en $2\pi$. El conjunto es permutacionalmente invariante, por lo que una arquitectura de tipo Deep Sets puede escribirse como
\begin{equation}
z_{\calP_m}=\rho\left(\sum_{\ell}\phi(\vect p_{m,\ell})\right).
\label{eq:deepset_paths}
\end{equation}
Esta forma admite un número variable de trayectorias y mantiene la correspondencia con el modelo físico; su invariancia a permutaciones sigue la construcción de Deep Sets \citep{Zaheer2017DeepSets}. Un transformador de conjuntos o un grafo pueden extenderla cuando las relaciones entre trayectorias sean relevantes. Para el alineamiento entre tecnologías de radio se aprenderá una representación común del soporte $\Gamma$, no una igualdad artificial entre los descriptores de $\calP_m$.

\section{Hipótesis de factorización latente}

Se propone separar
\begin{equation}
Z=\{Z_E,Z_\phi,Z_R,Z_H,Z_F,Z_N\},
\label{eq:factor_latent}
\end{equation}
con $Z_E$ asociado al entorno, $Z_\phi$ a la estructura e incertidumbre de fase, $Z_R$ a la configuración de radio, $Z_H$ al macroestado humano, $Z_F$ a la fisiología y $Z_N$ a perturbaciones residuales. La factorización no supone independencia absoluta; busca que la información relevante sea manipulable y que el codificador no confunda la identidad del dispositivo con el estado humano.

El bloque $Z_\phi$ puede dividirse operacionalmente en estructura predecible,
perturbaciones que deben estimarse o eliminarse y discrepancia residual, según la
\cref{eq:phase_latent_roles}. Esta clasificación no asigna una causa física única.
En particular, una pendiente espectral puede reunir retardo geométrico y error de
sincronización. Los diagnósticos propios que motivan la asignación de cada componente
se reservan al \cref{ch:avances}; aquí sólo se establece que la clasificación debe
revisarse cuando cambia la referencia de adquisición o la información auxiliar.

La identificabilidad de representaciones separadas requiere restricciones e
información adicional \citep{Locatello2019Identifiability}. Esta partición no es identificable
a partir de una única observación sin supuestos adicionales: distintas transformaciones invertibles de las variables latentes pueden explicar el mismo $Y_m$. En consecuencia, cada bloque tendrá criterios observables: intervenciones simuladas que modifican un factor a la vez, etiquetas parciales, pares de la misma escena observados con radios distintos, predicción cruzada y pruebas de invariancia fuera de distribución. La tesis evaluará esas propiedades y no interpretará automáticamente cada coordenada latente como una causa física.

\section{Dos ramas complementarias de representación}
\label{sec:two_representation_branches}

Se propone conservar simultáneamente una rama robusta $z_{\rm rob}=f_{\rm rob}(Y)$
y una rama coherente $z_{\rm coh}=f_{\rm coh}(Y,Y_{\rm ref},M)$, donde $M$ incluye
metadatos de adquisición y $Y_{\rm ref}$ una referencia independiente cuando exista.
La primera describe estructura relativa estable frente a un grupo instrumental
especificado; la segunda conserva variaciones temporales de fase y amplitud necesarias
para inferir movimiento. Su fusión se condicionará por la calidad de referencia y
la observabilidad, y se comparará con cada rama por separado.

El requisito deseado es invariancia a perturbaciones instrumentales e identificación
de la respuesta a intervenciones físicas. No puede garantizarse sólo mediante una
arquitectura: si una deriva instrumental y un movimiento generan la misma rotación
común, eliminarlos algebraicamente elimina ambas contribuciones. La separación exige
referencias, intervenciones controladas o información adicional. La equivariancia
física se expresará respecto de una transformación conocida del estado; no se
supondrá que toda traslación de un objeto produce una rotación global del canal.

\subsection{Familia de observables relativos}
\label{sec:representation_v_family}

La rama robusta puede construirse mediante observables con invariancias conocidas.
Para $H\in\mathbb C^{A\times K}$, $\mathbf h_k=H_{:,k}$ y $\epsilon>0$ con las
unidades del denominador, ejemplos representativos son
\begin{align}
\mathsf V_0(H)&=H, &\mathsf V_1(H)&=|H|,\nonumber\\
[\mathsf V_2^{(\ell)}(H)]_{a,k}&=\frac{H_{a,k+\ell}H_{a,k}^{*}}
{|H_{a,k+\ell}||H_{a,k}|+\epsilon},\qquad \ell\in\mathcal L,
\label{eq:v_frequency_products}\\
[\mathsf V_3(H)]_{a,b,k}&=\frac{H_{a,k}H_{b,k}^{*}}
{|H_{a,k}||H_{b,k}|+\epsilon},\label{eq:v_antenna_products}\\
\mathsf V_5(H)_k&=\frac{\mathbf h_k\mathbf h_k^{H}}
{\|\mathbf h_k\|_2^2+\epsilon},\label{eq:v_spatial_gram}\\
[\mathsf V_6(H)]_{a,k}&=H_{a,k}\exp(-j\widehat\psi+j2\pi\nu_k\widehat\tau).
\label{eq:v_canonical_channel}
\end{align}
$\mathsf V_4$ designa la comparación en el cociente por $U(1)$, mediante
$d_{U(1)}^2(x,y)=\|x\|^2+\|y\|^2-2|x^Hy|$, para el canal apilado de un arreglo.
Es una distancia entre clases de equivalencia, no un codificador unario. Los
productos espectrales conservan diferencias entre frecuencias y eliminan una fase
común; los productos espaciales eliminan fase común por frecuencia; el canal
canónico intenta estimar y retirar los parámetros en lugar de descartarlos. Estas
operaciones suponen grupos de perturbación diferentes y no deben agruparse bajo la
etiqueta genérica de fase robusta.

La selección de retardos espectrales, pares de antena, máscaras de desvanecimiento,
$\epsilon$ y normalización se realiza con entrenamiento y validación. Además,
$\widehat\psi$ y $\widehat\tau$ sólo son entradas legítimas en prueba si proceden
de una referencia disponible, metadatos o un predictor congelado. Alinear una
predicción con el canal objetivo define un límite oráculo. Los nombres de ejecución
y las comparaciones numéricas pertenecen al \cref{ch:avances}; la función de esta
sección es fijar el objeto matemático que allí se evalúa.

\section{Dinámica, inversión y transferencia}

WaveVerse amplía la simulación hacia mundos dinámicos y coherencia temporal
\citep{Zheng2026WaveVerse}; RayLoc e InverTwin utilizan diferenciabilidad para resolver
problemas inversos \citep{Han2026RayLoc,Chen2025InverTwin}; RFDT conecta una
representación RF con tareas posteriores \citep{Chen2026RFDT}. Estas capacidades no
son equivalentes. Una secuencia visualmente coherente no garantiza un Jacobiano correcto,
y una inversión exitosa no identifica parámetros físicos cuando el modelo directo es
no inyectivo.

La transferencia se evaluará con igual información, presupuesto de adaptación y
particiones espaciales. El preentrenamiento autosupervisado, incluida la familia JEPA
\citep{Assran2023IJEPA,Kim2026WiFiJEPA}, se considera un mecanismo para regularizar
la representación; su valor dependerá de preservar la respuesta a perturbaciones y de
reducir datos reales en dominios reservados.

Los insumos geométricos también forman parte de la comparación. DeepSense y CAEZ
aportan referencias de datos inalámbricos asociados con contexto y geometría
\citep{Alkhateeb2022DeepSense,ETHCAEZ,Zumegen2024CAEZWiFi}. WiSegRT condiciona
trazado y segmentación semántica \citep{Zhang2023WiSegRT}, mientras el mapeo interior
basado en radio y la identificación diferenciable de materiales investigan problemas
inversos sobre escena y propiedades \citep{Milosheski2026IndoorMapping,Chen2026RFMatID}.
WiInSim representa una vía industrial de simulación condicionada por escena
\citep{QualcommWiInSim}. Estos recursos no se compararán como si ofrecieran la misma
supervisión: se registrarán geometría, semántica, mediciones RF y parámetros materiales
disponibles en entrenamiento e inferencia.

\section{Comparación por familias de representación}

\begin{table}[htbp]
\centering\scriptsize
\caption{Familias de representación y función dentro de la tesis.}
\label{tab:rf_representation_comparison}
\begin{tabularx}{\textwidth}{P{3.2cm}P{2.0cm}L L}
\toprule
Familia & Estructura conservada & Ventaja buscada & Riesgo o contraste necesario\\
\midrule
Campo implícito \citep{Zhao2023NeRF2} & Continuidad espacial aprendida. &
Interpolación compacta del campo. & Ligazón a escena y datos; derivadas no
validadas por el error estático.\\
Geometría y objetos \citep{Jiang2025LearnableWDT,Chen2025RFScape,Zhang2026RadTwin}
& Escena, objetos y condición de radio. & Edición y transferencia estructurada. &
Más información de entrada; material y geometría pueden confundirse.\\
Trayectorias/conjuntos \citep{Zaheer2017DeepSets} & Retardos, ángulos,
interacciones y simetría de permutación. & Interfaz física de dimensión variable. &
Caminos no resolubles y correspondencias discontinuas.\\
Latente autosupervisada \citep{Assran2023IJEPA,Kim2026WiFiJEPA} & Regularidades
definidas por la tarea de preentrenamiento. & Menor supervisión y adaptación. &
El pretexto puede descartar sensibilidad física.\\
Relativa/invariante & Relaciones espaciales o espectrales. & Rechazo de
perturbaciones especificadas. & La dirección eliminada puede contener movimiento.\\
Coherentemente referenciada & Fase y amplitud respecto de una referencia. &
Micromovimiento e interpretación diferencial. & Dependencia de sincronización,
deriva y disponibilidad de referencia.\\
\bottomrule
\end{tabularx}
\end{table}

Ninguna familia domina por definición. La comparación debe igualar información,
presupuesto de adaptación y particiones, y debe medir tanto error estático como
respuesta a perturbaciones. El capítulo siguiente examina qué tareas humanas han
sido demostradas por Wi-Fi y LoRa; el \cref{ch:fidelity_observability} convierte el
requisito de preservar sensibilidad en un criterio matemático.
